\subsection{\texorpdfstring{\(p\)-根式扩张}{p-根式扩张}}

定义2.——设E是域K的一个扩域。我们称E是p-根式的 \(^{1}\) （在K上）如果E的每个元素在K上都是p-根式的。当这种情况发生时，若E的元素的高度集合有上界，则称E具有有限高度，并将E的高度定义为其元素高度的最大值。

我们注意到完全域（特别是特征为 0 的域）的每一个 \(p\) -根式扩张都是平凡的。

设K是一个域。K的高度为0的p-根式扩张是平凡扩张。K的每一个p-根式扩张都是代数的。若E是K的p-根式扩张，且F是E的p-根式扩张，则F是K的p-根式扩张：因为，给定任意 \(x \in F\) ，存在一个整数 \(m \geqslant 0\) 使得 \(x^{p^{m}} \in E\) 以及一个整数 \(n \geqslant 0\) 使得 \((x^{p^{m}})^{p^{n}} \in \mathbf{K}\) ，即 \(x^{p^{m+n}} \in K\) 。

命题2. --- 设E是域K的一个扩域。对于每个整数 \(n \geqslant 0\) ，令 \(E_{n}\) 为 E 中高度 \(\leqslant n\) （在 K 上）的 p-根式元素之集合，且设 \(E_{\infty}\) 为 E 中所有在 K 上为 p-根式元素的集合。则 \((\mathrm{E}_{n})_{n \geqslant 0}\) 是 E 的子扩张的一个递增序列，其并集为 \(E_{\infty}\) ，且 \(E_{\infty}\) 是包含在 E 中的 K 的最大 p-根式扩张。

对于每个整数 \(n \geqslant 0\) ，集合 \(E_{n}\) 由 E 中所有满足 \(x^{p^{n}} \in K\) 的元素 x 组成；由于映射 \(x \mapsto x^{p^{n}}\) 是域 E 的一个自同态，我们推断出 \(E_{n}\) 是 E 的一个子扩张。序列 \((\mathrm{E}_{n})_{n \geqslant 0}\) 是递增的，其并集为 \(E_{\infty}\) ，因此 \(E_{\infty}\) 是 E 的一个子扩张（V，第 11 页，命题 3）。显然， \(E_{\infty}\) 是 K 的一个 p-根式扩张，并且 \(E_{\infty}\) 包含 E 的每一个在 K 上是 p-根式的子扩张。

推论。--- 若域K的扩域E由K上的一组p-根式元素生成，则E在K上是p-根式的。

因为 \(E_{\infty}\) 是 E 的一个子扩张，且根据假设 \(\mathrm{E} = \mathrm{K}(\mathrm{E}_{\infty})\) ，由此 \(E = E_{\infty}\) ，所以 E 是 K 的 p-根式扩张。

在命题2的条件下， \(E_{\infty}\) 被称为 K 在 E 中的相对 p-根式闭包。

我们特别将命题2应用于E是K的代数闭扩张的情形；此时E是完全域，并且我们有 \(E_{n} = K^{p^{-n}}\) 对所有 \(n \geqslant 0\) 。在这种情况下，我们用 \(K^{p^{-\infty}}\) 表示E中在K上为p-根式的元素集合；它是E的子域，是E的子域的递增序列 \((\mathbf{K}^{p^{-n}})_{n \geqslant 0}\) 的并集。根据命题2， \(K^{p^{-\infty}}\) 是E在K上为p-根式的最大的子扩张；根据V第6页命题3， \(K^{p^{-\infty}}\) 是K的一个完全闭包，并且它也是E中包含K的最小的完全子域。当K是完全域时，我们显然有 \(K = K^{p^{-n}} = K^{p^{-\infty}}\) 对所有 n。若 K 是非完全域，我们有 \(K \neq K^{p}\) ，从而

\[
\mathbf {K} ^ {p ^ {- n}} \neq (\mathbf {K} ^ {p}) ^ {p ^ {- n}} = \mathbf {K} ^ {p ^ {- (n - 1)}}
\]

给定 \(n \geqslant 1\) ；子域 \(\mathbf{K}^{p^{-n}}\) （\(E\) 的）于是两两互不相同，并且 \(\mathbf{K}^{p^{-\infty}}\) 是 \(K\) 的一个无限次代数扩张。

命题3. --- 设 K 是一个域，E 是 K 的一个在 K 上为 p-根式的扩域，且 u 是 K 到完全域 F 内的一个同态。则存在唯一的将 u 扩张的 E 到 F 内的同态 v。

设 \(E_n\) 为 \(E\) 中 \(p\) -根式、高度 \(\leqslant n\) （在 \(K\) 上）的元素之集。根据命题2，域 \(E\) 是子域的递增序列 \((E_n)_{n \geqslant 0}\) 的并集。设 \(v\) 是 \(E\) 到 \(F\) 内且扩张 \(u\) 的一个同态；对任意 \(x \in E_n\) 我们有 \(x^{p^n} \in K\) ，从而 \(v(x)^{p^n} = v(x^{p^n}) = u(x^{p^n})\) ；于是我们有 \(v(x) = u(x^{p^n})^{p^{-n}}\) （对所有 \(n \geqslant 0\) 及所有 \(x \in E_n\)），由此得到 \(u\) 到 \(E\) 上的扩张的唯一性。

设 \(n\) 是一个正整数；对任意 \(x \in \mathrm{E}_n\) 我们有 \(x^{p^n} \in \mathbf{K}\) ，且由于 \(F\) 是完全的，我们可以定义一个元素 \(v_n(x)\) （属于 \(F\)）如下：\(v_n(x) = u(x^{p^n})^{p^{-n}}\) 。显然 \(v_n\) 是 \(E_n\) 到 \(F\) 内的一个同态，\(v_0 = u\) ，且 \(v_{n+1}\) 诱导 \(v_n\) （在 \(E_n\) 上）。于是存在一个同态 \(v\) （从 \(E\) 到 \(F\) 内），它诱导 \(v_n\) （在 \(E_n\) 上）对所有 \(n \geqslant 0\) ，特别地，与 \(v_0 = u\) 在 \(E_0 = K\) 上相一致。

推论。--- 要使域 K 的扩张 E 是 K 的完全闭包，必要且充分的条件是：它是 K 的 p-根式扩张，且域 E 是完全的。

当 p = 1 时，该推论是平凡的；现假设 \(p \neq 1\) 。由 V，第 5 页，定理 3 知所述条件是必要的；由命题 3 知它们是充分的。

命题4. --- 设E是域K的有限次p-根式扩张。则 {[}E : K{]} 是K的特征指数p的幂。

由于E是K的有限次p-根式扩张，存在 E 的元素 \(a_{1}, \ldots, a_{m}\) （在 K 上是 p-根式的），使得 \(\mathrm{E} = \mathrm{K}(a_{1}, \ldots, a_{m})\) 。设 i 在 1 到 m 的范围内；因为 \(a_{i}\) 关于 \(\mathrm{K}(a_{1}, \ldots, a_{i-1})\) 更是 p-根式的，次数

\[
n _ {i} = \left[ \mathrm{K} (a _ {1}, \dots , a _ {i}): \mathrm{K} (a _ {1}, \dots , a _ {i - 1}) \right]
\]

是 \(p\) 的幂（V，第 24 页，命题 1）。由 V，第 18 页，定理 2，我们有 \([\mathrm{E}:\mathbf{K}] = n_1\ldots n_m\) ，由此得结果。

